\section{Summary}
\label{sec_chapter2_summary}

This chapter has reviewed eleven papers directly relevant to autonomous, LLM-driven penetration testing, covering early single-agent systems (Fang et al., PentestGPT), multi-agent orchestration frameworks (HPTSA, VulnBot), planning-augmented approaches (CHECKMATE), non-web attack surfaces (Incalmo, PrediQL), and the two benchmark studies that provide the quantitative evidence base for the field's current limitations (CVE-Bench, PentestEval). The comparative analysis in Section~\ref{sec_comparative_analysis} identified five structural patterns across this literature: exploration-capable systems lack vulnerability-level dependency modeling; dependency-modeling systems require pre-enumerated weakness sets and do not operate on dynamically discovered output; cross-session memory is absent from every reviewed system; the two benchmark studies converge on the same compound failure from opposite ends; and no autonomous system in the reviewed set has been evaluated across more than one attack-surface family.

Section~\ref{sec_research_gap} translates these patterns into a precise three-dimensional gap: no reviewed system combines open-ended exploration of an unfamiliar attack surface with a dynamically constructed prerequisite-aware dependency model and evaluation across independently benchmarked attack-surface families. This gap is the problem this thesis investigates. Chapter~\ref{chapter_proposedMethod} introduces the proposed RedGrid architecture --- the four-layer orchestration hierarchy, Dual-Layer World Model, and VDG --- as a first design response to this gap. The architecture is at an early stage of specification and will receive further justification as implementation proceeds.
